\section{Limitaciones}
\label{sec:limitaciones}

El trabajo presenta las siguientes limitaciones. Primero, la distancia se
calcula con Haversine sobre una esfera de radio 6\,371\,000 m y se evalúa
contra los vértices del recorrido, no contra los segmentos; el backend
aplica el mismo criterio, por lo que el modelo y el código son
consistentes, pero la métrica es una aproximación de la distancia
peatonal y puede sobreestimar la cercanía en segmentos largos. Segundo,
por compatibilidad con la versión de la herramienta, los tipos registro
se representaron como tipos producto (tuplas), lo que reduce la
legibilidad del modelo frente a una definición con campos con nombre.
Tercero, la validación se basó en interpretación y casos de prueba, no en
una demostración deductiva completa de las obligaciones de prueba.
Finalmente, el conjunto de datos de prueba es reducido y no cubre toda la
variabilidad de las rutas reales.

Adicionalmente, el modelo NuSMV es una abstracción del ciclo de la
solicitud: no modela el algoritmo ni los datos de las rutas, y el backend
actual es stateless. Por tanto, la verificación con NuSMV no sustituye la
validación del código Django/React ni la validación del algoritmo en
VDM++.
